\section*{Chapitre \ref{ch:g10:reaction-progress-table} --- Le tableau d'avancement}

\begin{solution}{exo:g10:reaction-progress-table:1}
\ce{H2}~: $4.0 - 2x$~; \ce{O2}~: $3.0 - x$~; \ce{H2O}~: $2x$. Le
dihydrogène est épuisé pour $x = 2.0$, le dioxygène pour $x = 3.0$~:
$x_{\max} = \qty{2.0}{mol}$, le dihydrogène est limitant. \omterm{def:g10:reaction-progress-table:system}{État final}~:
\qty{0}{mol} de \ce{H2}, \qty{1.0}{mol} de \ce{O2}, \qty{4.0}{mol} de
\ce{H2O}.
\end{solution}

\begin{solution}{exo:g10:reaction-progress-table:2}
Le carbone serait épuisé pour $x = 3.0$, le dioxygène pour $x = 2.0$~:
$x_{\max} = \qty{2.0}{mol}$, le dioxygène est limitant. \omterm{def:g10:reaction-progress-table:system}{État final}~:
\qty{1.0}{mol} de carbone, pas de dioxygène, \qty{2.0}{mol} de dioxyde
de carbone.
\end{solution}

\begin{solution}{exo:g10:reaction-progress-table:3}
$x_{\max} = \qty{0.30}{mol}$ (le soufre est limitant). \omterm{def:g10:reaction-progress-table:system}{État final}~:
\qty{0.20}{mol} de fer, pas de soufre, \qty{0.30}{mol} de sulfure de
fer.
\end{solution}

\begin{solution}{exo:g10:reaction-progress-table:4}
Le diazote serait épuisé pour $x = 1.0$, le dihydrogène pour
$x = 2.4/3 = 0.80$~: $x_{\max} = \qty{0.80}{mol}$. \omterm{def:g10:reaction-progress-table:system}{État final}~:
\qty{0.20}{mol} de \ce{N2}, pas de \ce{H2},
$2 \times 0.80 = \qty{1.6}{mol}$ de \ce{NH3}.
\end{solution}

\begin{solution}{exo:g10:reaction-progress-table:5}
(b)~: $4/2 = 2/1$. En (a), le monoxyde de carbone est limitant~; en
(c) aussi.
\end{solution}

\begin{solution}{exo:g10:reaction-progress-table:6}
$n(\ce{CH4}) = 32 / 16.0 = \qty{2.0}{mol}$, donc $x_{\max} = \qty{2.0}{mol}$
pour un \omterm{def:g10:reaction-progress-table:stoichiometric}{mélange stœchiométrique}. Dioxygène~: $2 \times 2.0 = \qty{4.0}{mol}$,
soit $4.0 \times 32.0 = \qty{128}{g}$. Eau~:
$\qty{4.0}{mol} \times \qty{18.0}{g/mol} = \qty{72}{g}$.
\end{solution}

\begin{solution}{exo:g10:reaction-progress-table:7}
$n(\ce{C3H8}) = 1000 / 44.0 = \qty{22.7}{mol}$~; dioxyde de carbone~:
$3 \times 22.7 = \qty{68.2}{mol}$~; $V = 68.2 \times 24.1 =
\qty{1.64e3}{L}$, environ \qty{1.6}{m^3}.
\end{solution}

\begin{solution}{exo:g10:reaction-progress-table:8}
Pour $x = \qty{1.0}{mol}$~: \qty{1.0}{mol} de \ce{CH4}, \qty{1.0}{mol}
de \ce{O2}, \qty{1.0}{mol} de \ce{CO2}, \qty{2.0}{mol} de \ce{H2O}. Le
dioxygène est épuisé pour $x = \qty{1.5}{mol}$~: la réaction s'arrête
là, et les \omterm{def:g10:reaction-progress-table:extent}{avancements} plus grands ne sont jamais atteints.
\end{solution}

\begin{solution}{exo:g10:reaction-progress-table:9}
$n_0(\ce{Zn}) = 1.30 / 65.4 = \qty{0.0199}{mol}$~;
$n_0(\ce{Cu^{2+}}) = 0.10 \times 0.1000 = \qty{0.0100}{mol}$. Les \omterm{def:g9:ions:ion}{ions}
cuivre sont limitants~: $x_{\max} = \qty{0.0100}{mol}$. Cuivre formé~:
$0.0100 \times 63.5 = \qty{0.635}{g}$~; zinc restant~:
$(0.0199 - 0.0100) \times 65.4 = \qty{0.65}{g}$. Il ne reste plus d'\omterm{def:g9:ions:ion}{ions}
cuivre~: la couleur bleue a disparu.
\end{solution}

\begin{solution}{exo:g10:reaction-progress-table:10}
$n_0(\ce{Mg}) = 0.48 / 24.3 = \qty{0.0198}{mol}$, qui serait épuisé pour
$x = 0.0198$~; $n_0(\ce{H+}) = \qty{0.020}{mol}$, qui serait épuisé pour
$x = 0.010$. L'acide est limitant~: $x_{\max} = \qty{0.010}{mol}$, donc
\qty{0.010}{mol} de dihydrogène, $0.010 \times 24.1 = \qty{0.24}{L}$.
\end{solution}

\begin{solution}{exo:g10:reaction-progress-table:11}
Un \omterm{def:g10:reaction-progress-table:stoichiometric}{mélange stœchiométrique} avec \qty{0.30}{mol} de soufre demande
\qty{0.30}{mol} de fer~; il y en a \qty{0.20}{mol} de plus, soit
$0.20 / 0.30 \approx \qty{67}{\percent}$ d'excès.
\end{solution}

\begin{solution}{exo:g10:reaction-progress-table:12}
Quantités égales~: $m(\ce{Fe})/55.8 = m(\ce{S})/32.1$, avec
$m(\ce{Fe}) + m(\ce{S}) = \qty{10.0}{g}$. Donc
$m(\ce{Fe}) = 10.0 \times 55.8 / (55.8 + 32.1) = \qty{6.35}{g}$ et
$m(\ce{S}) = \qty{3.65}{g}$.
\end{solution}

\begin{solution}{exo:g10:reaction-progress-table:13}
\begin{enumerate}
  \item $n_0(\ce{CaCO3}) = 10.0 / 100.1 = \qty{0.0999}{mol}$, seul
        \omterm{def:g7:chemical-reactions:reactant}{réactif}~: $x_{\max} = \qty{0.0999}{mol}$~; on obtient
        \qty{0.0999}{mol} de \ce{CaO} et de \ce{CO2}, soit
        $0.0999 \times 24.1 = \qty{2.41}{L}$ de dioxyde de carbone.
  \item $n_0(\ce{CaO}) = \qty{0.0999}{mol}$, $n_0(\ce{H2O}) = 1.8 / 18.0
        = \qty{0.100}{mol}$~: la chaux vive est (tout juste) limitante~;
        le \omterm{def:g6:pure-substances-and-mixtures:mixture}{mélange} est presque stœchiométrique. \ce{Ca(OH)2}~:
        $M = 40.1 + 2 \times (16.0 + 1.0) = \qty{74.1}{g/mol}$, masse
        $0.0999 \times 74.1 = \qty{7.40}{g}$.
\end{enumerate}
\end{solution}

\begin{solution}{exo:g10:reaction-progress-table:14}
$n_0(\ce{H+}) = \qty{0.050}{mol}$ dans chaque erlenmeyer~; ces \omterm{def:g9:ions:ion}{ions}
seraient épuisés pour $x = \qty{0.025}{mol}$. Magnésium~: 0.15, 0.30,
0.45, 0.60, 0.75, \qty{0.90}{g} donnent 0.0062, 0.0123, 0.0185, 0.0247,
0.0309, \qty{0.0370}{mol}. Dans les quatre premiers erlenmeyers, le
magnésium est limitant et le dihydrogène vaut $n(\ce{Mg}) \times 24.1$~:
0.15, 0.30, 0.45 et \qty{0.60}{L}. Dans les deux derniers, l'acide est
limitant~: $0.025 \times 24.1
= \qty{0.60}{L}$. Le graphique est une droite passant par l'origine
jusqu'à environ \qty{0.61}{g} de magnésium, puis un palier horizontal à
\qty{0.60}{L}.
\end{solution}

\begin{solution}{exo:g10:reaction-progress-table:15}
$n_0(\text{éthanol}) = 4.6 / 46.0 = \qty{0.100}{mol}$ (épuisé pour
$x = 0.100$)~; $n_0(\ce{O2}) = 4.8 / 24.1 = \qty{0.199}{mol}$ (épuisé
pour $x = 0.199/3 = 0.0664$). Le dioxygène est limitant,
$x_{\max} = \qty{0.0664}{mol}$. Éthanol restant~: $0.100 - 0.0664 =
\qty{0.0336}{mol}$, soit $0.0336 \times 46.0 = \qty{1.5}{g}$. Dioxyde de
carbone~: $2 \times 0.0664 = \qty{0.133}{mol}$, soit
$0.133 \times 24.1 = \qty{3.2}{L}$.
\end{solution}

\begin{solution}{pb:g10:reaction-progress-table:1}
\textbf{1.} À gauche~: 2 Na, 6 N. À droite~: 2 Na, $3 \times 2 = 6$ N.
L'équation est ajustée.

\textbf{2.} À gauche~: 10 Na, 2 K, 2 N, 6 O. À droite~: K~: 2~; Na~:
$5 \times 2 =
10$~; O~: $1 + 5 = 6$~; N~: 2. L'équation est ajustée.

\textbf{3.} Le chapitre sur le \omterm{def:g9:periodic-table-first-look:periodic-table}{tableau périodique}
(\cref{ch:g9:periodic-table-first-look})~: le sodium réagit violemment
avec l'eau, en donnant une solution corrosive et du dihydrogène. Après
l'accident, le coussin est au contact de la peau, des yeux et de l'\omterm{def:g4:air-a-mixture-of-gases:air}{air}
humide de la respiration.

\textbf{4.} $M(\ce{NaN3}) = 23.0 + 3 \times 14.0 = \qty{65.0}{g/mol}$~;
$M(\ce{KNO3}) = 39.1 + 14.0 + 3 \times 16.0 = \qty{101.1}{g/mol}$.

\textbf{5.} \ce{NaN3}~: $1.00 - 2x$~; \ce{Na}~: $2x$~; \ce{N2}~: $3x$.

\textbf{6.} $1.00 - 2x = 0$ pour $x_{\max} = \qty{0.500}{mol}$~;
l'azoture de sodium, seul \omterm{def:g7:chemical-reactions:reactant}{réactif}, est limitant.

\textbf{7.} \ce{Na}~: $2 \times 0.500 = \qty{1.00}{mol}$~; \ce{N2}~:
$3 \times 0.500 = \qty{1.50}{mol}$.

\textbf{8.} $1.50 \times 24.1 = \qty{36.2}{L}$.

\textbf{9.} \ce{Na}~: $1.00 - 10x$~; \ce{KNO3}~: $n - 2x$~; \ce{K2O}~: $x$~;
\ce{Na2O}~: $5x$~; \ce{N2}~: $x$.

\textbf{10.} $1.00/10 = n/2$, donc $n = \qty{0.200}{mol}$, soit
$0.200 \times 101.1 = \qty{20.2}{g}$.

\textbf{11.} $x_{\max} = \qty{0.100}{mol}$~: pas de sodium, pas de nitrate
de potassium, \qty{0.100}{mol} de \ce{K2O}, \qty{0.500}{mol} de
\ce{Na2O}, \qty{0.100}{mol} de \ce{N2}.

\textbf{12.} Le sodium serait épuisé pour $x = 0.100$, le nitrate de
potassium pour $x = 0.150/2 = 0.075$~: le nitrate est limitant,
$x_{\max} =
\qty{0.075}{mol}$, et il resterait dans le coussin
$1.00 - 10 \times 0.075 = \qty{0.25}{mol}$ de sodium métallique, prêt à
réagir avec l'humidité.

\textbf{13.} $1.50 + 0.100 = \qty{1.60}{mol}$ de diazote par \omterm{def:g10:the-mole:amount}{mole}
d'azoture de sodium.

\textbf{14.} $n = 60 / 24.1 = \qty{2.49}{mol}$.

\textbf{15.} $2.49 / 1.60 = \qty{1.56}{mol}$ d'azoture de sodium.

\textbf{16.} $1.56 \times 65.0 = \qty{101}{g}$.

\textbf{17.} $0.200 \times 1.56 = \qty{0.311}{mol}$, soit
$0.311 \times 101.1 = \qty{31.5}{g}$ de nitrate de potassium.

\textbf{18.} $2.49 / 1.50 = \qty{1.66}{mol}$, soit
$1.66 \times 65.0 = \qty{108}{g}$~: la seconde réaction fait économiser
environ \qty{7}{g} d'azoture de sodium.

\textbf{19.} Une \omterm{def:g10:reaction-progress-table:limiting}{réaction totale} ne s'arrête que lorsque son \omterm{def:g10:reaction-progress-table:limiting}{réactif
limitant} est épuisé~; ici, l'azoture de sodium est le seul \omterm{def:g7:chemical-reactions:reactant}{réactif}, donc
dans l'\omterm{def:g10:reaction-progress-table:system}{état final} sa quantité vaut $1.00 - 2x_{\max} = 0$. Il ne reste
rien de toxique dans le coussin.

\textbf{20.} Environ \qty{101}{g} d'azoture de sodium remplissent un
airbag de \qty{60}{L}.
\end{solution}
